
\section{Hidden Variables and Bell's Proof}
\label{sec:Hidden Variables}

To arrive at Bell’s inequality, it is necessary to understand the concept of hidden variables and how it was used throughout the decades between the EPR paper and John Bell’s 1964 work. To trace the discussion back to its origins, it is important to identify when and by whom the idea was first introduced. This requires examining the work of von Neumann and understanding his impossibility theorem.

\subsection{Von Neumann's Axioms and the Impossibility of Hidden Variables}

In 1932, John von Neumann published \textit{Mathematische Grundlagen der Quantenmechanik}, establishing the first rigorous and axiomatic formulation of Quantum Mechanics. His goal was to determine whether the quantum‐mechanical description of reality was \textbf{complete} or if additional ``hidden'' parameters might restore determinism. 

This formal synthesis was possible thanks to the earlier demonstration of \textbf{compatibility between Schrödinger’s wave mechanics and Heisenberg’s matrix mechanics}. Schrödinger had shown that both formulations were mathematically equivalent representations of the same physical theory. Von Neumann formalized this equivalence as an \textbf{isomorphism built upon Hilbert spaces}, where physical observables are represented by self-adjoint operators and states by vectors or density operators within the same unified mathematical structure.

\vspace{0.3cm}

\subsubsection*{Axioms for the Expectation Value}

Von Neumann did not begin from the quantum postulates themselves, but from general properties that any consistent statistical theory of physical magnitudes should satisfy.  
He denoted by $\mathcal{V}_m(R)$ the mean value of a physical quantity (observable) $R$ in an ensemble, and introduced two fundamental axioms:

\begin{description}
    \item[A'. Positivity:] If $R$ is a non-negative magnitude (for instance, the square of another quantity), then
    \begin{equation}
        \mathcal{V}_m(R) \ge 0.
    \end{equation}

    \item[B'. Linearity:] For arbitrary magnitudes $R,S$ and real coefficients $a,b$,
    \begin{equation}
        \mathcal{V}_m(aR + bS) = a\,\mathcal{V}_m(R) + b\,\mathcal{V}_m(S),
    \end{equation}
    even when $R$ and $S$ do not commute. 
\end{description}

Under these assumptions, von Neumann proved that there exists a unique Hermitian, positive operator $U$ (today called the \textit{statistical operator} or density matrix) such that
\begin{equation}
    \mathcal{V}_m(R) = \mathrm{Tr}(UR).
\end{equation}

\vspace{0.3cm}

\subsubsection*{Conditions Required by a Hidden-Variable Theory}

If the quantum state were only a statistical mixture of more fundamental ``dispersion-free'' states characterized by hidden variables $\lambda$, each observable $R$ would take a definite value $v(R,\lambda)$ and
\begin{equation}
    \mathcal{V}_m(R) = \int v(R,\lambda)\,\rho(\lambda)\,d\lambda.
\end{equation}

Such a model must satisfy two additional conditions:

\begin{enumerate}
    \item \textbf{Eigenvalue condition:}
    \[
    \mathcal{V}_m(R) \in \{\text{eigenvalues of } R\}.
    \]
    \item \textbf{Dispersion-free condition:}
    \[
    \mathcal{V}_m(R^2) = [\mathcal{V}_m(R)]^2.
    \]
\end{enumerate}

These express that every individual system possesses precise values for all observables, and that quantum indeterminacy reflects only ignorance of $\lambda$.

\vspace{0.3cm}

\subsubsection*{The Incompatibility Proof}

Von Neumann demonstrated that no expectation-value function $\mathcal{V}_m(\cdot)$ can simultaneously satisfy both the linearity postulate (B′) and the dispersion-free condition for all observables.

If $R$ and $S$ are non-commuting operators, linearity would require
\begin{equation}
    v(R + S,\lambda) = v(R,\lambda) + v(S,\lambda),
\end{equation}
whereas the eigenvalue condition demands that $v(R + S,\lambda)$ be one of the eigenvalues of $R + S$, which in general is not the sum of eigenvalues of $R$ and $S$ taken separately.  
Therefore, for non-commuting observables $[R,S]\neq 0$, no assignment $v(R,\lambda)$ can fulfill all axioms simultaneously.

\begin{equation}
    \mathcal{V}_m(RS) + \mathcal{V}_m(SR) = 2\,\mathcal{V}_m(R)\,\mathcal{V}_m(S)
    \quad \text{cannot hold in general.}
\end{equation}

This contradiction constitutes von Neumann’s \textbf{Impossibility Theorem for Hidden Variables}:  
\begin{center}
\emph{No hidden-variable theory can reproduce all statistical predictions of quantum mechanics while preserving linearity and positivity of expectation values.}
\end{center}

\subsection{Hidden Variables by Bohm} 

Only three years after Einstein, Podolsky, and Rosen proposed their famous EPR argument questioning the completeness of Quantum Mechanics, it is interesting to consider whether Einstein was unaware of von Neumann’s 1932 work on hidden variables, or whether he deliberately ignored it to challenge the prevailing interpretation.  
Regardless of that, in 1952 David Bohm published two seminal papers that reopened the discussion.  
His work provided a concrete deterministic model based on \textit{hidden variables}, offering both a mathematical framework and a potential test of the EPR argument.

\subsubsection{Historical Context and Motivation}

In 1952, David Bohm introduced a deterministic and causal interpretation of Quantum Mechanics.  
He demonstrated that it was possible, in principle, to reproduce all the statistical predictions of quantum theory without abandoning realism or determinism.  
His formulation, directly inspired by the Einstein–Podolsky–Rosen argument, extended Louis de Broglie’s earlier ``pilot-wave'' concept into a fully consistent dynamical theory.

Bohm’s proposal appeared to contradict von Neumann’s 1932 \textit{Impossibility Theorem}, which claimed that no hidden-variable theory could reproduce the predictions of Quantum Mechanics.  
However, Bohm pointed out that von Neumann’s proof relied on an assumption that was \textbf{not physically necessary}: the \textit{linearity of expectation values for non-commuting observables}.  
By rejecting this assumption, Bohm’s theory avoided the theorem’s restriction and reopened the question of completeness.  
His interpretation restored determinism but at the cost of introducing explicit \textbf{nonlocal interactions}.

\subsubsection{The Pilot-Wave Theory}

Bohm’s reformulation begins with the non-relativistic Schrödinger equation for a complex wave function $\Psi(x,t)$, expressed in polar form as
\[
\Psi(x,t) = R(x,t)e^{iS(x,t)/\hbar},
\]
where $R$ and $S$ are real functions representing the amplitude and phase, respectively.  
Substituting this into the Schrödinger equation and separating real and imaginary parts yields two coupled equations with distinct physical meanings.

\paragraph{(1) The Continuity Equation.}
The imaginary part gives:
\begin{equation}
\frac{\partial R^2}{\partial t} + \nabla \cdot \left( R^2 \frac{\nabla S}{m} \right) = 0,
\end{equation}
which has the form of a continuity equation, with $R^2 = |\Psi|^2$ interpreted as a probability density and
\[
\vec{v} = \frac{\nabla S}{m}
\]
as the velocity field.  
This leads to the \textbf{Guidance Equation}, where each particle follows a definite trajectory determined by the phase $S$ of the pilot wave.

\paragraph{(2) The Modified Hamilton–Jacobi Equation.}
The real part gives:
\begin{equation}
\frac{\partial S}{\partial t} + \frac{(\nabla S)^2}{2m} + V(x) + U(x) = 0,
\end{equation}
where
\begin{equation}
U(x) = -\frac{\hbar^2}{2m}\frac{\nabla^2 R}{R}
\end{equation}
is the \textbf{quantum potential}.  
This term, absent in classical mechanics, depends on the spatial form of the wave rather than its magnitude.  
It represents a nonclassical influence responsible for interference and tunneling phenomena.  
In the double-slit experiment, for instance, $U(x)$ depends on the geometry of both slits, guiding each particle along deterministic trajectories toward regions of constructive interference—even though it passes through only one slit.

\subsubsection{Statistical Predictions and the Born Rule}

A central question in Bohm’s interpretation concerns the recovery of the probabilistic results of standard Quantum Mechanics.  
Although each particle follows a deterministic path, their initial positions are unknown.  
Their statistical distribution is governed by the \textbf{Born rule}, which in standard quantum theory assigns probability density
\begin{equation}
\rho(x,t) = |\Psi(x,t)|^2.
\end{equation}
Bohm introduced this as a \textbf{quantum equilibrium condition}: if the initial distribution of positions satisfies $\rho(x,0) = |\Psi(x,0)|^2$, then this relation remains valid for all subsequent times.  
Combining the guidance equation
\[
\frac{dx}{dt} = \frac{\nabla S}{m}
\]
with the continuity equation,
\[
\frac{\partial \rho}{\partial t} + \nabla \cdot (\rho\, \mathbf{v}) = 0,
\quad \text{where} \quad \mathbf{v} = \frac{\nabla S}{m},
\]
shows that both $\rho$ and $|\Psi|^2$ obey identical evolution laws.  
Therefore, if the system begins in quantum equilibrium, the ensemble will always reproduce the predictions of standard Quantum Mechanics:
\[
P(a \in [a_1,a_2]) = \int_{a_1}^{a_2} |\Psi(a,t)|^2\, da.
\]

The apparent randomness of measurement outcomes thus arises not from indeterminism, but from ignorance of the exact initial conditions.  
Consequently, Bohm’s deterministic theory yields \textbf{exactly the same statistical predictions} as the Copenhagen interpretation.

\subsubsection{Nonlocality and the EPR Problem}

The most striking feature of Bohm’s interpretation emerges when it is extended to systems with multiple particles.  
For two particles, the wave function and the quantum potential are defined in a \textbf{six-dimensional configuration space}:
\begin{equation}
\Psi(\mathbf{x}_1,\mathbf{x}_2,t)
= R(\mathbf{x}_1,\mathbf{x}_2,t)\,
  e^{iS(\mathbf{x}_1,\mathbf{x}_2,t)/\hbar}.
\end{equation}
The corresponding \textbf{quantum potential} is
\begin{equation}
U(\mathbf{x}_1,\mathbf{x}_2,t)
= -\sum_{i=1}^{2}
\frac{\hbar^2}{2m_i}
\frac{\nabla_i^2 R(\mathbf{x}_1,\mathbf{x}_2,t)}
     {R(\mathbf{x}_1,\mathbf{x}_2,t)}.
\end{equation}
This potential couples the motion of both particles through its dependence on their joint configuration $(\mathbf{x}_1,\mathbf{x}_2)$.  
A change in the position of one particle instantly alters the potential acting on the other, regardless of the distance separating them.

The corresponding \textbf{quantum force} acting on each particle is
\begin{equation}
\mathbf{F}_i = -\nabla_i (V + U)
= -\nabla_i V
+ \frac{\hbar^2}{2m_i}
\nabla_i \!\left( \frac{\nabla_i^2 R}{R} \right).
\end{equation}
Because $R$ depends on both coordinates, the force $\mathbf{F}_i$ is inherently \textbf{nonlocal}: the dynamics of one particle depend instantaneously on the configuration of the entire system.

This is the explicit manifestation of nonlocality in Bohm’s theory.  
In the EPR scenario, a measurement on one particle alters the total wave function $\Psi$, and through the quantum potential, this change instantaneously affects the motion of the distant particle.  
Bohm’s model thus provides a deterministic and causal account of the EPR correlations while remaining fully compatible with all statistical predictions of Quantum Mechanics.

\subsubsection{The Central Consequence and Bohm’s Legacy}

Bohm’s work demonstrated that a \textbf{deterministic hidden-variable theory} can reproduce all statistical predictions of Quantum Mechanics, contradicting the widely held interpretation of von Neumann’s theorem as an absolute no-go result.  
The key price to pay, however, is explicit nonlocality: quantum influences act instantaneously between distant systems.

This led to a profound dichotomy:
\[
\textbf{Determinism} \quad \Longleftrightarrow \quad \textbf{Nonlocality.}
\]
Any completion of quantum mechanics that restores determinism must necessarily abandon locality.  

This conceptual tension directly inspired John S. Bell’s work in 1964.  
Bell analyzed Bohm’s theory and proved that any hidden-variable model reproducing the predictions of Quantum Mechanics must violate local causality.  
He formalized this requirement as an experimentally testable constraint—the \textbf{Bell inequalities}.  
Their subsequent violation in experiments confirmed that \textbf{Bohm’s nonlocal hidden-variable theory and standard Quantum Mechanics make identical empirical predictions}, revealing that nonlocal correlations are indeed a fundamental feature of nature.


\subsection{John Bell}

The theoretical framework developed by David Bohm in 1952 showed that a deterministic hidden-variable theory could reproduce all statistical predictions of Quantum Mechanics, but only at the cost of introducing explicit nonlocality.  
More than a decade later, John S. Bell sought to determine whether such nonlocal features were unavoidable.  
In his 1966 paper \textit{“On the Problem of Hidden Variables in Quantum Mechanics”}, he critically analyzed von Neumann’s impossibility theorem and demonstrated that its central assumption—the linearity of expectation values—was not physically necessary.  
This opened the door to hidden-variable models like Bohm’s.  

Two years earlier, however, in 1964, Bell had already published his most celebrated result: \textit{“On the Einstein Podolsky Rosen Paradox.”}  
In this work, he derived a simple but powerful mathematical constraint—known as the \textbf{Bell inequality}—showing that the predictions of Quantum Mechanics are fundamentally incompatible with any theory that respects both \textbf{locality} and \textbf{realism}.  
Bell’s theorem thus transformed a philosophical question into a concrete, experimentally testable statement.

\subsubsection{Fundamental Assumptions: Realism and Locality}

Bell’s analysis starts from two physical assumptions that define the class of \textit{local hidden-variable theories}:

\begin{itemize}
    \item \textbf{Realism:}  
    Physical quantities have definite values prior to and independent of measurement.  
    The outcome of a spin measurement is determined by a set of hidden parameters $\lambda$, which fully describe the physical state of the system.

    \item \textbf{Locality:}  
    Measurement outcomes at one location are unaffected by measurement settings or results at distant, space-like separated regions.  
    No instantaneous influences are allowed between remote systems.
\end{itemize}

Under these principles, the result of a spin measurement on particle 1 along direction $\mathbf{a}$ is represented by a function $A(\mathbf{a},\lambda) = \pm 1$, and similarly for particle 2 along direction $\mathbf{b}$, $B(\mathbf{b},\lambda) = \pm 1$.  
The hidden variables $\lambda$ are distributed according to a normalized probability density $\rho(\lambda)$ characterizing the particle source.

\subsubsection{Local Correlations and Expectation Values}

The average correlation between results measured along directions $\mathbf{a}$ and $\mathbf{b}$ is given by:
\begin{equation}
P(\mathbf{a},\mathbf{b}) = \int d\lambda\, \rho(\lambda)\, A(\mathbf{a},\lambda)\, B(\mathbf{b},\lambda).
\label{eq:bell-correlation}
\end{equation}
Locality requires that $A$ depends only on $\mathbf{a}$ and $\lambda$, and $B$ only on $\mathbf{b}$ and $\lambda$.  
This structure embodies the idea that each result is locally determined, while statistical correlations arise solely from the shared hidden variables $\lambda$.


\subsubsection{Derivation and Interpretation of Bell’s Inequality}

Bell considered the measurement of spin components of two entangled particles along different directions $\mathbf{a}$, $\mathbf{b}$, and $\mathbf{c}$.  
Each measurement outcome is dichotomic, taking the values $+1$ or $-1$, corresponding to ``spin up'' or ``spin down'' relative to the chosen axis.

Under the assumptions of \textbf{realism} and \textbf{locality}, the outcomes are determined by hidden variables $\lambda$ according to:
\[
A(\mathbf{a},\lambda), \; A(\mathbf{b},\lambda), \; A(\mathbf{c},\lambda) \in \{-1, +1\}.
\]
For particle pairs in the singlet state, the results for both particles are perfectly anti-correlated when measured along the same direction:
\[
A(\mathbf{a},\lambda) = -B(\mathbf{a},\lambda).
\]
Using this relation, the correlation between measurements along $\mathbf{a}$ and $\mathbf{b}$ is written as:
\[
P(\mathbf{a},\mathbf{b}) = -\int d\lambda\, \rho(\lambda)\, A(\mathbf{a},\lambda)A(\mathbf{b},\lambda),
\]
where $\rho(\lambda)$ is the probability distribution of the hidden variables.

Now consider three different orientations of the measuring devices, $\mathbf{a}$, $\mathbf{b}$, and $\mathbf{c}$.  
The difference between two correlations can be expressed as:
\[
P(\mathbf{a},\mathbf{b}) - P(\mathbf{a},\mathbf{c})
= -\int d\lambda\, \rho(\lambda)\, A(\mathbf{a},\lambda)\left[A(\mathbf{b},\lambda) - A(\mathbf{c},\lambda)\right].
\]

Because $A(\mathbf{b},\lambda)$ and $A(\mathbf{c},\lambda)$ can only take values $\pm 1$, one can show algebraically that
\[
|A(\mathbf{b},\lambda) - A(\mathbf{c},\lambda)| = 1 - A(\mathbf{b},\lambda)A(\mathbf{c},\lambda),
\]
which leads directly to the inequality:
\[
|P(\mathbf{a},\mathbf{b}) - P(\mathbf{a},\mathbf{c})| \le 1 + P(\mathbf{b},\mathbf{c}).
\]
This is the \textbf{Bell inequality} in its original form.

\vspace{0.4em}
\noindent
\textbf{Physical interpretation of each term:}
\begin{itemize}
    \item The quantities $P(\mathbf{a},\mathbf{b})$, $P(\mathbf{a},\mathbf{c})$, and $P(\mathbf{b},\mathbf{c})$ represent the \textbf{correlation coefficients} between pairs of spin measurements along different directions.  
    Each $P(\mathbf{x},\mathbf{y})$ ranges from $-1$ (perfect anticorrelation) to $+1$ (perfect correlation).
    
    \item The left-hand side, $|P(\mathbf{a},\mathbf{b}) - P(\mathbf{a},\mathbf{c})|$, quantifies the \textbf{difference between two measurable correlations}.  
    It expresses how the correlation between measurements along $\mathbf{a}$ and $\mathbf{b}$ differs from that along $\mathbf{a}$ and $\mathbf{c}$, when the hidden variable $\lambda$ determines all outcomes consistently.
    
    \item The right-hand side, $1 + P(\mathbf{b},\mathbf{c})$, acts as a \textbf{local-realistic upper bound}.  
    The constant $1$ originates from the discrete nature of outcomes ($\pm 1$), while $P(\mathbf{b},\mathbf{c})$ captures the mutual correlation between the two directions $\mathbf{b}$ and $\mathbf{c}$ used for the comparison.  
    Together, they set the maximum variation in correlations allowed by any model where results are predetermined and unaffected by distant settings.
\end{itemize}

\vspace{0.4em}
\noindent
Thus, Bell’s inequality expresses a precise mathematical limit on how strongly measurement results at distant locations can be correlated, assuming that those results are fixed by local hidden variables.  
Any violation of this bound therefore signals the failure of the joint assumptions of locality and realism.

\subsubsection{Violation Predicted by Quantum Mechanics}

Quantum Mechanics predicts that for two spin-$\tfrac{1}{2}$ particles in the singlet state,
\[
P(\mathbf{a},\mathbf{b}) = -\mathbf{a}\cdot\mathbf{b} = -\cos(\theta_{ab}),
\]
where $\theta_{ab}$ is the angle between the measurement directions.  
Choosing, for instance,
\[
\mathbf{a} = 0^\circ, \quad \mathbf{b} = 45^\circ, \quad \mathbf{c} = 90^\circ,
\]
we obtain:
\[
P(\mathbf{a},\mathbf{b}) = -\tfrac{1}{\sqrt{2}}, \quad
P(\mathbf{a},\mathbf{c}) = 0, \quad
P(\mathbf{b},\mathbf{c}) = -\tfrac{1}{\sqrt{2}}.
\]
Substituting into Bell’s inequality yields:
\[
|P(\mathbf{a},\mathbf{b}) - P(\mathbf{a},\mathbf{c})| = \tfrac{1}{\sqrt{2}} > 1 + P(\mathbf{b},\mathbf{c}) = 0.293,
\]
showing a clear violation of the inequality.  
Thus, the predictions of Quantum Mechanics exceed the limits imposed by any local hidden-variable theory.

\subsubsection{Physical Implications}

Bell’s result establishes that no physical theory can reproduce all quantum predictions while maintaining both locality and realism.  
This finding marks a turning point in the foundations of Quantum Mechanics:

\begin{itemize}
    \item \textbf{Incompatibility:}  
    Local hidden-variable theories cannot account for the correlations predicted by Quantum Mechanics.

    \item \textbf{Testability:}  
    Bell transformed the EPR debate from a philosophical discussion into an experimentally testable proposition.

    \item \textbf{The Dilemma:}  
    Any successful description must abandon at least one classical notion:
    \[
    \textbf{Determinism with Locality} \;\;\not\Rightarrow\;\; \text{Quantum Predictions.}
    \]
\end{itemize}

Bell’s theorem thus revealed that if quantum predictions are correct—as experiments later confirmed—then the physical world is inherently \textbf{nonlocal} or \textbf{nonrealistic}.  
This result laid the foundation for all subsequent experimental tests of entanglement and the modern development of quantum information theory.

\subsubsection{Bell’s Response to Bohm and von Neumann (1966)}

In 1966, John S. Bell published a less-known but foundational paper titled \textit{“On the Problem of Hidden Variables in Quantum Mechanics.”}  
This work was motivated directly by David Bohm’s 1952 model, which had demonstrated the logical possibility of a deterministic hidden-variable theory reproducing all quantum predictions.  
Bell’s goal was to clarify how Bohm’s construction could exist despite von Neumann’s earlier ``proof'' of impossibility.

Von Neumann’s 1932 theorem had concluded that no hidden-variable theory could reproduce the statistical results of Quantum Mechanics.  
His reasoning relied on two assumptions regarding the structure of observables and expectation values:

\begin{enumerate}
    \item \textbf{Linearity of Expectation Values:}
    \[
    \langle A + B \rangle = \langle A \rangle + \langle B \rangle,
    \]
    assumed to hold for all observables $A$ and $B$, even when they do not commute.
    
    \item \textbf{Positivity of Values:}  
    The expected value of a positive operator must always be non-negative, ensuring that $\langle \psi | A^\dagger A | \psi \rangle \ge 0$.
\end{enumerate}

Bell argued that while these assumptions are valid within standard Quantum Mechanics, they are not necessary constraints for alternative theories that assign definite values to observables.  
In particular, he showed that von Neumann’s linearity postulate implicitly assumes the quantum structure it seeks to exclude: it enforces additivity of expectation values even for non-commuting quantities, which cannot be measured simultaneously in any physical experiment.

By relaxing this assumption, Bell demonstrated that the impossibility theorem collapses, and that \textbf{hidden-variable models are not logically excluded}.  
In his words, von Neumann’s proof ``was not merely false but foolish''—it eliminated hidden variables only by imposing conditions they could never satisfy.

Bell emphasized that Bohm’s 1952 pilot-wave model stands as a counterexample: it assigns definite positions to particles while preserving all measurable predictions of Quantum Mechanics.  
What von Neumann had ruled out was not Bohm’s kind of theory, but only a class of linear expectation-value assignments incompatible with physical realism.

\begin{quote}
    ``The impossibility proof of von Neumann was not relevant to the question of whether there could be hidden variables in quantum mechanics.''  
    \hfill — \textit{J. S. Bell, 1966}
\end{quote}

The 1966 paper thus served two purposes: it vindicated Bohm’s construction as mathematically consistent, and it established the conceptual foundation for Bell’s later 1964 theorem by clarifying what assumptions must be questioned when seeking a deeper theory.  
In essence, Bell’s critique reopened the logical space for hidden-variable theories, but his own 1964 result showed that any such completion must inevitably be \textbf{nonlocal}.

\subsection{The CHSH Inequality (1969)}

While Bell’s 1964 theorem established the incompatibility between Quantum Mechanics and all local hidden-variable theories, its original inequality was not directly suited for laboratory tests.  
In 1969, John Clauser, Michael Horne, Abner Shimony, and Richard Holt reformulated Bell’s argument into a version that could be experimentally verified with entangled photons or spin pairs.  
This result, known as the \textbf{CHSH inequality}, became the standard framework for all subsequent experimental tests of quantum nonlocality.

\subsubsection{Reformulation of Bell’s Argument}

The CHSH scenario considers two observers, Alice and Bob, located in distant regions, each choosing between two possible measurement settings:
\[
A, A' \quad \text{for Alice,} \qquad B, B' \quad \text{for Bob.}
\]
Each measurement yields a dichotomic result $A, A', B, B' \in \{-1, +1\}$ determined by local settings and hidden variables $\lambda$, distributed according to $\rho(\lambda)$.

The correlation function for outcomes measured along settings $A$ and $B$ is defined as:
\begin{equation}
E(A,B) = \int d\lambda\, \rho(\lambda)\, A(A,\lambda)\, B(B,\lambda).
\end{equation}

Under the joint assumptions of \textbf{realism} and \textbf{locality}, the following combination must satisfy the inequality:
\begin{equation}
S = E(A,B) + E(A',B) + E(A,B') - E(A',B'),
\qquad
|S| \le 2.
\label{eq:CHSH}
\end{equation}

\noindent
\textbf{Physical meaning of each term:}
\begin{itemize}
    \item Each correlation $E(X,Y)$ measures the statistical dependence between outcomes observed by Alice and Bob under particular choices of detector orientations.  
    \item The combination $S$ captures the maximal contrast achievable between sets of four such correlations, assuming that the results for $A, A', B, B'$ are all predetermined by $\lambda$ and that no instantaneous influence exists between the two sides.
    \item The upper bound $|S|\le 2$ represents the strict limit imposed by any local hidden-variable model.  
    Any stronger correlation would require either nonlocal interactions or contextual dependence of outcomes.
\end{itemize}

\subsubsection{Quantum Mechanical Prediction}

Quantum Mechanics predicts that for a maximally entangled two-spin system (the singlet state), the correlation between spin measurements along directions $\mathbf{a}$ and $\mathbf{b}$ is:
\[
E(\mathbf{a},\mathbf{b}) = -\cos(\theta_{ab}),
\]
where $\theta_{ab}$ is the angle between the measurement settings.  
Substituting these correlations into Eq.~\eqref{eq:CHSH} with the optimal choice of relative angles,
\[
\theta_{AB} = 0^\circ, \quad
\theta_{A'B} = 45^\circ, \quad
\theta_{AB'} = 45^\circ, \quad
\theta_{A'B'} = 90^\circ,
\]
yields:
\[
S_{\text{QM}} = 2\sqrt{2} \approx 2.828,
\]
which clearly violates the local-realistic bound of $2$.  
This is the celebrated \textbf{quantum violation of the CHSH inequality}.

\subsubsection{Experimental Verification}

The CHSH formulation provided a practical way to test the foundations of Quantum Mechanics.  
Beginning with the pioneering experiment of Freedman and Clauser (1972), followed by Alain Aspect’s optical tests in 1982, and later the ``loophole-free'' experiments by Weihs et al. (1998) and Hensen et al. (2015), all results have confirmed the quantum mechanical prediction:
\[
|S_{\text{exp}}| > 2.
\]
No experiment consistent with relativistic separation and detector efficiency has ever satisfied the local-realistic bound.

\subsubsection{Consequences and Legacy}

The violation of the CHSH inequality has profound implications for our understanding of nature:

\begin{itemize}
    \item \textbf{Nonlocality as a feature of nature:}  
    Correlations between entangled particles exceed any possible classical explanation based on local realism.
    
    \item \textbf{From foundations to technology:}  
    The same quantum correlations responsible for Bell violations form the basis of quantum communication, teleportation, and cryptography.
    
    \item \textbf{Conceptual shift:}  
    The CHSH result transformed Bell’s theorem from a theoretical constraint into an empirical fact:  
    \[
    \text{Quantum Mechanics} \Rightarrow |S| > 2 \quad \Longrightarrow \quad \text{Nature is nonlocal.}
    \]
\end{itemize}

Thus, the CHSH inequality closed the logical circle begun by Einstein, Bohm, and Bell, showing that quantum nonlocality is not just a matter of interpretation but an experimentally verified property of the physical world.
